\SourceLeaf{309}%
das Verfahren anzugeben, mit dem sich diese Entwicklung beliebig weit fortführen lässt.

Ich erlaube mir einige Änderungen an den dort verwendeten Bezeichnungen, teils um den Druck zu erleichtern, teils um die Bezeichnungen in den einzelnen Teilen der vorliegenden Abhandlung etwas symmetrischer zu gestalten. Zunächst bedeuten hier
\begin{list}{}{\setlength{\leftmargin}{1.4em}\setlength{\itemindent}{-1.4em}\setlength{\itemsep}{0pt}\setlength{\parsep}{0pt}}
\item \(r\) den Abstand der beiden Punkte voneinander, wobei der Radius der Kugel als Einheit dient.
\item \(B+\frac12b\) und \(B-\frac12b\) die Breiten am ersten und zweiten Endpunkt von \(r\).
\item \(T+\frac12t\) und \(T-\frac12t\pm180^0\) die Azimute des zweiten Endpunkts vom ersten aus beziehungsweise des ersten Endpunkts vom zweiten aus.
\item \(l\) den Längenunterschied.
\end{list}
\noindent Es wird angenommen, dass das Azimut von Süden nach Westen gezählt wird und \(l\) positiv ist, wenn der zweite Punkt weiter westlich liegt als der erste.

Ferner setzen wir
\[
\begin{aligned}
\beta&=r\cos T\\
\tau&=r\sin T\mathord{.}\operatorname{tang}B\\
\lambda&=\frac{r\sin T}{\cos B}
\end{aligned}
\]
Diese Größen entsprechen dem, was in Artikel 17 mit \(b^0,a^0,\lambda^0\) bezeichnet wurde, nämlich den bis zur dritten Ordnung ausschließlich genauen Werten von \(b,t,l\). Zwischen ihnen gilt die Gleichung
\[
rr+\tau\tau=\beta\beta+\lambda\lambda
\]
Die Ordnungen sind hier stets so zu verstehen, dass \(r\) als Größe erster Ordnung betrachtet wird.

\indent Zur Abkürzung schreiben wir außerdem
\[
\begin{aligned}
\frac{2\operatorname{tang}\frac12r}{r}&=m\\
\frac{2\sin\frac12r}{r}&=n
\end{aligned}
\]
Zur beabsichtigten Entwicklung gelangt man am leichtesten, indem man die Gleichung
\[
x=\sin y
\]
